\section{引言}
\label{sec:intro}
部分子分布函数（PDF）编码了强相互作用重要的非微扰信息。基于 QCD 因子化~\cite{Collins:1989gx}，人们已经成功地以很好的精度从高能散射数据中抽取了 PDF~\cite{Lin:2017snn}。然而，无论从理论还是实用的角度来看，从第一性原理的格点 QCD（LQCD）计算中抽取 PDF 都必须去做：这既是为了检验 QCD 的非微扰区，也是为了研究那些难以进行散射实验的强子的部分子结构。

由于定义 PDF 的算符具有时间依赖性，在欧氏空间的 LQCD 中{\it 直接}计算 PDF 即使并非不可能，也是困难的~\cite{Lin:2017snn}。Ji 提出了一种新方法~\cite{Ji:2013dva}，他引入了一组可在 LQCD 中计算的准部分子分布（quasi-PDF），并论证了当强子动量 $P_z$ 被推进到无穷大时，准部分子分布趋于相应的部分子分布。从 LQCD 计算中抽取 PDF 的其他一些方法也被提出过~\cite{Liu:1993cv,Liu:1999ak,Liu:2016djw,Orginos:2017kos,Chambers:2017dov}。在文献~\cite{Ma:2014jla,Ma:2017pxb}中，我们中的两位作者提出了一种基于 QCD 因子化的普适方法，用以\textit{间接}地在 LQCD 中计算 PDF。
与从可因子化、可测量的强子截面实验数据中抽取 PDF 类似，我们提议通过对 LQCD 计算得到的良好“格点截面”（lattice cross sections，LCSs）数据进行全局分析来抽取 PDF；LCSs 定义为满足以下条件的强子矩阵元：(1) 可在欧氏空间 LQCD 中计算；(2) 对紫外（UV）发散可重正化，以保证可靠的连续极限；(3) 可通过红外安全的匹配系数因子化为 PDF。正是第 (3) 条因子化把所求的 PDF 与 LQCD 可计算的 LCSs 联系起来。

为了抽取丰富、精确且按味分离的 PDF 信息，必须找到尽可能多的良好 LCSs，因为不同味的 PDF 很可能对同一个 LQCD 计算的 LCS 有贡献。为构造良好的 LCSs，我们研究了两类算符：（a）带有适当规范链的两个规范依赖场算符的关联，我们称之为准部分子算符，因为它们涵盖了定义准部分子分布的所有算符乃至更多；（b）两个规范不变流算符的关联。在我们的方案中，在坐标空间中对每个良好 LCS 进行 LQCD 计算提供了\textit{约束} PDF 所需的信息，这与测量各种实验截面以\textit{约束} PDF 的作用类似。作为对比，在 Ji 的方案~\cite{Ji:2013dva}中，关注点在于用在足够大 $P_z$ 下 LQCD 计算的准部分子分布来重现相应的 PDF。



对于 LQCD 计算，与 (b) 类算符相比，计算 (a) 类算符的花费相对较低，而 (b) 类算符的重正化则比 (a) 类简单得多。(b) 类算符的重正化几乎是平凡的：只需重正化众所周知规范不变的定域流即可。另一方面，(a) 类算符的重正化并不平凡，原因在于相应算符的非定域性和可能的幂次发散。两类算符的 QCD 因子化已在文献~\cite{Ma:2017pxb}中研究过，我们发现这些算符的乘法可重正性是共线因子化成立的必要条件。

人们已投入大量努力来探究 (a) 类准夸克算符的紫外结构~\cite{Ishikawa:2016znu,Chen:2016fxx,Monahan:2016bvm,Briceno:2017cpo, Xiong:2017jtn, Constantinou:2017sej, Alexandrou:2017huk, Chen:2017mzz}。
准夸克算符的全阶乘法可重正性已用两种不同方法得到证明：一种依赖辅助场技术~\cite{Ji:2017oey,Green:2017xeu}，另一种基于图展开~\cite{Ishikawa:2017faj}。这些证明为“从 LQCD 计算的准夸克算符强子矩阵元中抽取对胶子不敏感的那类夸克分布组合”提供了坚实的理论基础~\cite{Lin:2014zya,Alexandrou:2015rja,Chen:2016utp,Alexandrou:2016jqi,Zhang:2017bzy,Monahan:2017hpu,Chen:2017lnm,Stewart:2017tvs,Izubuchi:2018srq,Alexandrou:2018pbm,Chen:2018xof,Chen:2018fwa,Liu:2018uuj,Bali:2018spj,Radyushkin:2018nbf,Lin:2018qky,Karpie:2018zaz}。

本文在 QCD 微扰论的所有阶研究准胶子算符的紫外发散。正如我们将要解释的，准胶子算符的紫外结构可能复杂得多。我们把一般裸准胶子算符定义为
\begin{align}\label{eq:ftgx}
\begin{split}
{\cal O}_{bg}^{\mu \nu \rho\sigma}(\xi) = F^{\mu \nu}(\xi)\,\Phi^{(a)}(\{\xi,0\})\, F^{{\rho\sigma}}(0)\,,
\end{split}
\end{align}
其中 $\Phi^{(a)}(\xi,0)={\cal P}e^{-ig_s\int_0^{1} \xi\cdot A^{(a)}(r \xi)\,dr}$ 是伴随表示下路径排序的规范链。为明确起见，我们设 $\xi_\mu$ 沿 $z$ 方向，并引入单位矢量 $n^\mu = (0,0,0, 1)$，对任意矢量 $v_\mu$ 定义 $v\cdot n \equiv v_z$。由于 $F^{\mu\nu}$ 中含有量纲导数算符，表观幂次计数告诉我们，胶子场强与规范链之间的顶点可能是线性紫外发散的。使用截断正规化时，文献~\cite{Wang:2017qyg,Wang:2017eel}中的一圈计算确实显示出该顶点存在未消去的线性发散，这会使准胶子算符的乘法重正化几乎不可能实现。相比之下，准夸克算符的相应顶点至多是对数紫外发散的。

另一个复杂性来自 UV 重正化中的算符混合。原则上，考虑到胶子场强的反对称性之后，式~\eqref{eq:ftgx} 中的一般准胶子算符可以有多达 36 个独立算符，它们在重正化下都可能相互混合。在文献中，准胶子 PDF 由 ${\cal O}_{bg}^{\mu \nu \rho\sigma}$ 与某些“张量”缩并后的线性组合构造而成。文献~\cite{Ji:2013dva}、\cite{Wang:2017eel}以及首个格点模拟~\cite{Fan:2018dxu}分别通过用 $n_\mu n_\rho  g_{\nu i}g_\sigma^{i}$、$n_\mu n_\rho  g_{\nu \sigma}$ 和 $g_{\mu 0}g_{\rho0}  g_{\nu \sigma}$ 缩并 ${\cal O}_{bg}^{\mu \nu \rho\sigma}$ 得到了不同的准胶子 PDF 定义。由于这些算符之间的混合，这些准胶子 PDF 的重正化是非平凡的。

本文首先对动量为 $p$ 的渐近胶子的准胶子算符进行显式一圈计算，其定义为 $\langle g(p)|{\cal O}_{bg}^{\mu \nu \rho\sigma}(\xi) | g(p) \rangle$。我们采用维数正规化（DR）同时正规化对数与线性紫外发散：在 $n$ 圈阶，它们分别表现为围绕 $d=4$ 和 $d=4-1/n$ 的极点。我们发现，在 DR 下胶子—规范链顶点的一圈修正中的线性紫外发散会相消，这使得准胶子算符的乘法可重正性成为可能。随后我们考察了高阶图所有可能的紫外发散拓扑。利用规范不变性，我们发现来自胶子—规范链顶点的全部线性紫外发散在微扰论所有阶都相消。
进而我们发现，全部 36 个独立准胶子算符都可以被乘法重正化而不与任何其他算符混合。结合我们在文献~\cite{Ishikawa:2017faj}中关于准夸克算符的证明，本文针对准胶子算符的工作完成了准部分子算符乘法重正化的证明。
